\documentclass[10pt,a4paper]{article}

% ============ Пакеты ============
\usepackage{amsmath,amssymb,amsthm}
\usepackage{mathtools}
\usepackage[margin=2.3cm]{geometry}   % поля симметричны: left=right=2.3cm
\usepackage{fontspec}
\usepackage[russian]{babel}
\usepackage{xcolor}
\usepackage{booktabs}
\usepackage{tabularx}
\usepackage{longtable}
\usepackage{enumitem}
\usepackage{needspace}
% hyperref -- последний среди контентных пакетов
\usepackage[
  colorlinks=true,
  linkcolor=linkblue,
  citecolor=darkgray,
  urlcolor=linkblue,
  bookmarks=true,
  bookmarksnumbered=true,
  unicode=true
]{hyperref}

% ============ Шрифты (кириллица, Times-метрика) ============
\setmainfont{Liberation Serif}
\setsansfont{Liberation Sans}

% ============ Цвета ============
\definecolor{linkblue}{RGB}{28,58,122}
\definecolor{darkgray}{RGB}{70,70,70}

% ============ Теоремные окружения ============
\theoremstyle{plain}
\newtheorem{theorem}{Теорема}[section]
\newtheorem{lemma}[theorem]{Лемма}
\newtheorem{proposition}[theorem]{Предложение}
\newtheorem{corollary}[theorem]{Следствие}
\theoremstyle{definition}
\newtheorem{definition}[theorem]{Определение}
\newtheorem{example}[theorem]{Пример}
\theoremstyle{remark}
\newtheorem{remark}[theorem]{Замечание}
\renewcommand{\proofname}{Доказательство}

% ============ Макросы ============
\newcommand{\Om}{\boldsymbol{\Omega}}
\newcommand{\om}{\boldsymbol{\omega}}
\newcommand{\R}{\mathbb{R}}
\newcommand{\so}{\mathfrak{so}(3)}
\newcommand{\RP}{\mathbb{RP}}
\newcommand{\bbI}{\mathbf{I}}
\newcommand{\bbL}{\mathbf{L}}
\newcommand{\bbP}{\mathbf{P}}
\newcommand{\bu}{\mathbf{u}}
\newcommand{\bvn}{\mathbf{v}}
\newcommand{\bwn}{\boldsymbol{\omega}}
\newcommand{\bx}{\mathbf{x}}
\newcommand{\bn}{\mathbf{n}}
\DeclareMathOperator{\spec}{spec}
\DeclareMathOperator{\Tr}{Tr}
\DeclareMathOperator{\adst}{ad^{*}}
\DeclareMathOperator{\diag}{diag}

% ============ Типографика ============
\widowpenalty=10000
\clubpenalty=10000
\tolerance=2500
\hbadness=3000
\setlength{\emergencystretch}{3em}
\numberwithin{equation}{section}

% ============ Метаданные ============
\hypersetup{
  pdftitle={Математическая модель P^3 -> so(3) для Навье--Стокса: строгое изложение, расширения и численная верификация},
  pdfauthor={POLER Research Notes},
  pdfsubject={Геометрическая редукция динамики завихренности на алгебру Ли so(3)},
  pdfkeywords={Навье--Стокс, so(3), Эйлер--Пуанкаре, Ли--Пуассон, Казимир, Хопф--Коул, МакВини, регулярность}
}

\begin{document}

% Титульная страница отсутствует в .tex: обложка генерируется отдельно (HTML/Playwright)
% и подклеивается как страница 0. Первая страница тела -- оглавление.
\tableofcontents
\newpage

\needspace{0.25\textheight}
\section*{Аннотация}
\addcontentsline{toc}{section}{Аннотация}

Работа даёт полное строгое изложение конечномерной модели динамики завихренности
Навье--Стокса на алгебре Ли $\so \cong \R^3$: геометрия проективного замыкания
$\RP^3 \cong \mathrm{SO}(3)$, вывод уравнений Эйлера--Пуанкаре с диссипацией,
инварианты Казимира и гамильтониан, каскад редукции $P^3 \to 2D \to 1D \to 0D$
(теорема Ладыженской, линеаризация Коула--Хопфа, спектральный проектор МакВини)
и главный результат об отсутствии сингулярностей в модели. Все утверждения
сопровождаются полными доказательствами; помимо исходной программы изложены
расширения: прямое доказательство тождества Якоби скобки Ли--Пуансона,
теорема периодичности орбит, оптимальность константы $4$ в оценке дефекта
идемпотентности МакВини и условная теорема о тензоре деформации, уточняющая
границы переноса результатов на полное уравнение Навье--Стокса. Раздел~8
содержит численную верификацию всех тождеств и оценок (23 проверки,
интегратор RK4): инварианты сохраняются с точностью $5{,}7\cdot 10^{-9}$
при порядке сходимости $p\approx 4{,}96$, орбиты замкнуты с точностью
$5\cdot 10^{-15}$, оценка Грёнуолла со штрейном достигается точно.

\bigskip

\needspace{0.25\textheight}
\section{Введение и постановка задачи}\label{sec:intro}

Пусть вязкая несжимаемая жидкость занимает область $\Omega \subset \R^3$,
поле скорости $\bu(x,t)$ удовлетворяет системе Навье--Стокса, а завихренность
$\bwn = \nabla \times \bu$~--- её вихревой форме
\begin{equation}\label{eq:ns-vort}
  \partial_t \bwn + (\bu \cdot \nabla)\bwn
  = (\bwn \cdot \nabla)\bu + \nu \Delta \bwn,
  \qquad \nabla \cdot \bu = 0 .
\end{equation}
Центральный вопрос глобальной регулярности \eqref{eq:ns-vort}~--- может ли норма
$\|\bwn(\cdot,t)\|_{L^\infty}$ достигнуть бесконечности за конечное время
$T^{*}$ \cite{fefferman}. Ключевую роль играет расщепление нелинейности:
обозначим $S = \tfrac12(\nabla\bu + (\nabla\bu)^{\mathsf T})$ тензор деформации
(симметричная часть $\nabla\bu$). Член растяжения вихря допускает тогда
точное представление $(\bwn\cdot\nabla)\bu = S\bwn$ (лемма~\ref{lem:strain}):
вихрь не «вращается», а только деформируется симметричной частью поля.
Именно этот член способен увеличивать $\|\bwn\|$, и именно он делает
трёхмерную проблему открытой.

\emph{Программа настоящей работы} (модель $P^3 \to \so$) состоит в замене
бесконечномерной динамики \eqref{eq:ns-vort} конечномерной системой
Эйлера--Пуанкаре на $\so^{*} \cong \R^3$ для амплитудного вектора
вихревой трубки $\Om = (\Omega_1,\Omega_2,\Omega_3)$:
\begin{equation}\label{eq:epd-intro}
  \dot{\Om} = \Om \times (\bbI^{-1}\Om) - \nu\,\bbL\,\Om,
\end{equation}
где $\bbI = \diag(I_1,I_2,I_3)$, $I_k > 0$~--- тензор инерции трубки,
$\bbL = \bbL^{\mathsf T} \succeq 0$~--- оператор вязкого затухания
(структурный аналог $-\Delta$), а нелинейность
$\Om \mapsto \Om \times (\bbI^{-1}\Om)$~--- чисто кососимметричный
поворот фазы, ортогональный $\Om$ при всех $t$. Модель сохраняет
геометрию задачи: состояние живёт на компактном многообразии
$\RP^3 \cong \mathrm{SO}(3)$, а $\Om$~--- в коалгебре $\so^{*}$.

Содержание работы~--- строгий разбор модели от геометрии до границ
применимости. Перечислим основные результаты; каждый пункт снабжён
ссылкой на точную формулировку.

\begin{enumerate}[label=\textup{(\roman*)}, leftmargin=2.2em, itemsep=2pt]
  \item \textbf{Геометрия} (\S\ref{sec:geom}): полный вывод диффеоморфизма
        $\RP^3 \cong \mathrm{SO}(3)$ через кватернионы, включая явную
        формулу поворота ось--угол (теорема~\ref{th:rp3so3}); тождества
        алгебры $\so \cong \R^3$ с метрикой Киллинга
        $\langle u,v\rangle = -\tfrac12\Tr(\hat u \hat v) = u\cdot v$
        (леммы~\ref{lem:hatjacobi}, \ref{lem:killing}).
  \item \textbf{Вывод системы} (\S\ref{sec:ep}): вариационный вывод
        \eqref{eq:epd-intro} из редуцированного лагранжиана
        (лемма~\ref{lem:variation}, теорема~\ref{th:ep}); скобка
        Ли--Пуансона $\{F,G\} = -\langle\Om, \nabla F \times \nabla G\rangle$
        с \emph{полным прямым доказательством} тождества Якоби
        (предложение~\ref{prop:poisson}); глобальная корректность задачи
        Коши (теорема~\ref{th:wp}).
  \item \textbf{Инварианты} (\S\ref{sec:inv}): точные законы эволюции
        Казимира $C = \|\Om\|^2$ и гамильтониана
        $\mathcal H = \tfrac12\langle\Om,\bbI^{-1}\Om\rangle$
        (теоремы~\ref{th:casimir}, \ref{th:ham}); экспоненциальное
        затухание $C(t) \le C(0)e^{-2\nu\lambda_{\min}t}$
        (следствие~\ref{cor:decay}); периодичность всех неравновесных
        орбит при $\nu = 0$ (предложение~\ref{prop:periodic});
        классификация и устойчивость равновесий, включая эффект
        «теннисной ракетки» (предложение~\ref{prop:stability}).
  \item \textbf{Каскад редукции} (\S\ref{sec:cascade}):
        точная лемма о растяжении $(\bwn\cdot\nabla)\bu = S\bwn$
        (лемма~\ref{lem:strain}); теорема Ладыженской об энстрофии
        (теорема~\ref{th:ladyzh}); полная линеаризация Коула--Хопфа
        (теорема~\ref{th:colehopf}); спектральная теория проектора
        МакВини $\mathcal Q(\bbP) = 3\bbP^2 - 2\bbP^3$ с доказательством
        \emph{оптимальности константы} $4$ в оценке дефекта
        (теорема~\ref{th:mcweeny}, предложение~\ref{prop:mcweeny-iter}).
  \item \textbf{Отсутствие сингулярностей} (\S\ref{sec:noblowup}):
        главный результат~--- априорная оценка
        $\sup_{t \ge 0}\|\Om(t)\| \le \|\Om(0)\| < \infty$ при всех
        $\nu \ge 0$ (теорема~\ref{th:main}); явные осесимметричные
        решения с вязким замедлением прецессии (предложение~\ref{prop:axisym}).
  \item \textbf{Границы метода} (\S\ref{sec:critique}): точное структурное
        объяснение, почему модель не решает трёхмерную проблему
        Клэя (несферичность коалгебраических орбит группы
        диффеоморфизмов против сферических орбит $\so$), и условная
        теорема со штрейном $\dot\Om = \ldots + S(t)\Om - \nu\bbL\Om$,
        воспроизводящая структуру критериев BKM/CKN
        (теорема~\ref{th:strain}).
  \item \textbf{Верификация} (\S\ref{sec:numerics}): численная проверка
        всех тождеств и оценок: 23/23 тестов, включая машинно-точную
        замкнутость орбит и точное достижение границы Грёнуолла.
\end{enumerate}

Мы сознательно отделяем доказанное внутри модели от вопроса о переносе
на уравнение \eqref{eq:ns-vort}: модельная теорема~\ref{th:main} полна и
строга, однако её перенос на трёхмерное уравнение требует контроля
$\|S(t)\|$, который модель постулирует, а полное уравнение не даёт;
подробное обсуждение~--- в \S\ref{sec:critique}. Нумерация всех величин
символьна: параметры $I_1, I_2, I_3$, $\nu$, $\bbL$ нигде не фиксируются,
кроме численных экспериментов \S\ref{sec:numerics}.

\needspace{0.25\textheight}
\section{Геометрия многообразия: кватернионы, \texorpdfstring{$\RP^3 \cong \mathrm{SO}(3)$}{RP3=SO(3)} и алгебра \texorpdfstring{$\so$}{so(3)}}\label{sec:geom}

\subsection{Кватернионы и двойное накрытие}

\begin{definition}[алгебра кватернионов]\label{def:quat}
$\mathbb H$~--- четырёхмерная вещественная алгебра с базисом
$1, e_1, e_2, e_3$ и таблицей умножения
$e_1e_2 = e_3$, $e_2e_3 = e_1$, $e_3e_1 = e_2$,
$e_ie_j = -e_jei$ при $i \ne j$, $e_i^2 = -1$. Сопряжение:
$\overline{q} = a - be_1 - ce_2 - de_3$ для $q = a + be_1 + ce_2 + de_3$;
норма $|q|^2 = q\bar q = a^2 + b^2 + c^2 + d^2$. Подпространство чисто
мнимых кватернионов обозначим $\operatorname{Im}\mathbb H \cong \R^3$.
\end{definition}

Множество единичных кватернионов $S^3 = \{q : |q| = 1\}$~--- группа
по умножению, поскольку $|q_1q_2| = |q_1||q_2|$. Зададим отображение
\begin{equation}\label{eq:conj-action}
  R(q)v := q\,v\,\bar q, \qquad q \in S^3,\; v \in \operatorname{Im}\mathbb H .
\end{equation}

\begin{lemma}\label{lem:conjugation}
Отображение \eqref{eq:conj-action} корректно определено на
$\operatorname{Im}\mathbb H$, является гомоморфизмом
$R(q_1q_2) = R(q_1)R(q_2)$ и принимает значения в $\mathrm{SO}(3)$.
\end{lemma}

\begin{proof}
Чистота: $\overline{qv\bar q} = q\bar v \bar q = -qv\bar q$ при
$\bar v = -v$, поэтому $qv\bar q \in \operatorname{Im}\mathbb H$.
Изометрия: $|qv\bar q| = |q|\,|v|\,|\bar q| = |v|$, значит $R(q)$~---
ортогональное линейное отображение $\R^3 \to \R^3$. Гомоморфизм:
$R(q_1q_2)v = q_1q_2 v \bar q_2 \bar q_1 = R(q_1)(R(q_2)v)$.
Определитель: $S^3$ связно, $\det R(\cdot)$ непрерывен и принимает
значения $\pm 1$, причём $\det R(1) = 1$; следовательно
$\det R(q) \equiv 1$ и $R(q) \in \mathrm{SO}(3)$.
\end{proof}

\begin{lemma}[формула ось--угол]\label{lem:axisangle}
Пусть $q = \cos\tfrac\theta2 + e_3\sin\tfrac\theta2$. Тогда в базисе
$e_1, e_2, e_3$
\begin{equation}\label{eq:rodrigues}
  R(q)e_1 = \cos\theta\, e_1 + \sin\theta\, e_2,\qquad
  R(q)e_2 = -\sin\theta\, e_1 + \cos\theta\, e_2,\qquad
  R(q)e_3 = e_3 .
\end{equation}
\end{lemma}

\begin{proof}
Прямое умножение с учётом $e_3e_1 = e_2$, $e_1e_3 = -e_2$,
$e_2e_3 = e_1$, $e_3e_2 = -e_1$, $e_i^2 = -1$; положим
$c = \cos\tfrac\theta2$, $s = \sin\tfrac\theta2$. Тогда
\begin{align}
  q e_1 \bar q &= (ce_1 + se_3e_1)(c - se_3)
   = (ce_1 + se_2)(c - se_3)
   = c^2e_1 - cse_1e_3 + cse_2 - s^2 e_2e_3 \notag\\
   &= (c^2 - s^2)e_1 + 2cs\,e_2 = \cos\theta\,e_1 + \sin\theta\,e_2,
   \label{eq:e1}\\
  q e_2 \bar q &= (ce_2 - se_1)(c - se_3)
   = c^2e_2 - cse_2e_3 - cse_1 + s^2e_1e_3
   = \cos\theta\,e_2 - \sin\theta\,e_1, \label{eq:e2}\\
  q e_3 \bar q &= (ce_3 - s)(c - se_3)
   = c^2e_3 + cs - cs + s^2e_3 = (c^2+s^2)e_3 = e_3. \label{eq:e3}
\end{align}
Формулы \eqref{eq:e1}--\eqref{eq:e3} и есть поворот на угол $\theta$
вокруг оси $e_3$.
\end{proof}

\begin{theorem}[$\RP^3 \cong \mathrm{SO}(3)$]\label{th:rp3so3}
Отображение $R\colon S^3 \to \mathrm{SO}(3)$~--- сюръективный гомоморфизм
групп с ядром $\{\pm 1\}$. Оно индуцирует гомеоморфизм (и диффеоморфизм)
\[
  S^3/\{v \sim -v\} \;=\; \RP^3 \;\longrightarrow\; \mathrm{SO}(3),
\]
в частности $\mathrm{SO}(3)$ компактно.
\end{theorem}

\begin{proof}
\emph{Ядро.} Пусть $R(q) = \mathrm{Id}$, т.е. $qv = vq$ для всех
$v \in \operatorname{Im}\mathbb H$. Расписывая $q = a + be_1 + ce_2 + de_3$
и приравнивая $qe_1 = e_1q$, получаем $ce_3 - de_2 = -ce_3 + de_2$, т.е.
$c = d = 0$; аналогично из $qe_2 = e_2q$ получаем $b = d = 0$.
Значит $q = a \in \R \cap S^3 = \{\pm1\}$.

\emph{Сюръективность.} Пусть $A \in \mathrm{SO}(3)$. Характеристический
многочлен $A$ веществен и нечётной степени, $\det A = 1$, все собственные
значения по модулю равны $1$; комплексные корни идут парами
$\lambda, \bar\lambda$, и перемножая, получаем, что $1$~--- собственное
значение: если все корни вещественны, то $\{1,1,1\}$ либо
$\{1,-1,-1\}$; если есть пара $e^{\pm i\alpha}$, то третий корень равен
$1$. Пусть $\bn$~--- соответствующий единичный собственный вектор.
Ограничение $A$ на плоскость $\bn^\perp$~--- элемент $\mathrm{SO}(2)$,
поворот на некоторый угол $\theta$. Выбирая в $\bn^\perp$
ортонормированный базис и сопрягая, сводим задачу к случаю
$\bn = e_3$, разобранному в лемме~\ref{lem:axisangle}:
$A = R(\cos\tfrac\theta2 + \bn\sin\tfrac\theta2)$ (формулы
\eqref{eq:rodrigues}~--- это в точности поворот на $\theta$ вокруг $\bn$
в подходящем базисе; смена базиса поглощается выбором знака $\theta$).

\emph{Фактор.} Индуцированное отображение $\bar R\colon S^3/\{\pm1\} \to
\mathrm{SO}(3)$~--- непрерывная биекция. Факторпространство $S^3/\{\pm1\}$
компактно как образ компакта при непрерывной проекции, $\mathrm{SO}(3)$
хаусдорфово; непрерывная биекция из компакта в хаусдорфово пространство~---
гомеоморфизм \cite{marsdenratiu}. Гладкость $\bar R$ и обратного к нему~--- стандартный факт (биективная иммерсия компактного многообразия);
см.~\cite{marsdenratiu}. По определению $S^3/\{v\sim -v\} = \RP^3$.
\end{proof}

\subsection{Алгебра \texorpdfstring{$\so$}{so(3)} и метрика Киллинга}

\begin{definition}[шляп-изоморфизм]\label{def:hat}
Для $u \in \R^3$ положим
\[
  \hat u := \begin{pmatrix} 0 & -u_3 & u_2\\ u_3 & 0 & -u_1\\ -u_2 & u_1 & 0 \end{pmatrix}
  \in \so .
\]
Прямая проверка даёт $\hat u\, v = u \times v$ для всех $v \in \R^3$.
\end{definition}

\begin{lemma}[тождество Якоби для $\times$]\label{lem:jacobi-cross}
Для всех $x, y, z \in \R^3$
\[
  x \times (y \times z) + y \times (z \times x) + z \times (x \times y) = 0 .
\]
\end{lemma}

\begin{proof}
Формула Лагранжа $x\times(y\times z) = y\,(x\cdot z) - z\,(x\cdot y)$
даёт
\[
  \sum_{\text{цикл}} x\times(y\times z)
  = \bigl[y(x\cdot z) - z(x\cdot y)\bigr]
  + \bigl[z(y\cdot x) - x(y\cdot z)\bigr]
  + \bigl[x(z\cdot y) - y(z\cdot x)\bigr] = 0,
\]
так как скалярное произведение коммутативно и каждое слагаемое
$w\,(v\cdot u)$ встречается дважды с противоположными знаками.
\end{proof}

\begin{lemma}\label{lem:hatjacobi}
Изоморфизм $u \mapsto \hat u$ удовлетворяет
$\,[\hat u, \hat v] = \widehat{u \times v}$; следовательно
$(\R^3, \times)$~--- алгебра Ли, изоморфная $\so$.
\end{lemma}

\begin{proof}
Для любого $w \in \R^3$, пользуясь определением~\ref{def:hat} и
леммой~\ref{lem:jacobi-cross},
\[
  [\hat u,\hat v]\,w = \hat u(v\times w) - \hat v(u \times w)
  = u\times(v\times w) - v\times(u\times w) = (u\times v)\times w,
\]
где последний переход~--- перестановка тождества Якоби:
$u\times(v\times w) - v\times(u\times w)
= -w\times(u\times v) = (u\times v)\times w$.
Совпадение линейных отображений на всех $w$ даёт
$[\hat u,\hat v] = \widehat{u\times v}$.
\end{proof}

\begin{lemma}[метрика Киллинга]\label{lem:killing}
Форма
$\langle u, v\rangle := -\tfrac12\Tr(\hat u\hat v)$
удовлетворяет
\[
  \langle u,v\rangle = u\cdot v \quad\text{и}\quad
  \langle [u,v], w\rangle = \langle u, [v,w]\rangle
  \quad\text{(ad-инвариантность).}
\]
\end{lemma}

\begin{proof}
Запишем $\hat u_{ij} = -\varepsilon_{ijk}u_k$. Тогда
\[
  \Tr(\hat u\hat v)
  = \sum_{i,j} \hat u_{ij}\hat v_{ji}
  = \sum_{i,j,k,\ell} \varepsilon_{ijk}\varepsilon_{j i \ell}\,
    u_k v_\ell
  = -\sum_{k,\ell} u_kv_\ell \sum_{i,j}\varepsilon_{ijk}\varepsilon_{ij\ell}
  = -2\sum_{k}u_kv_k = -2\,u\cdot v,
\]
где использовано $\sum_{i,j}\varepsilon_{ijk}\varepsilon_{ij\ell}
= 2\delta_{k\ell}$. Отсюда $\langle u,v\rangle = u\cdot v > 0$ при
$u \ne 0$. Ad-инвариантность есть цикличность смешанного произведения:
\[
  \langle u\times v, w\rangle = (u\times v)\cdot w
  = u\cdot(v\times w) = \langle u, v\times w\rangle. \qedhere
\]
\end{proof}

Положительная определённость $\langle\cdot,\cdot\rangle$ означает, что
$\so$~--- алгебра Ли компактного типа; это отражает компактность
$\mathrm{SO}(3)$ из теоремы~\ref{th:rp3so3}.

\begin{definition}[тензор инерции вихревой трубки]\label{def:inertia}
Симметричный положительно определённый оператор
$\bbI = \diag(I_1, I_2, I_3)$, $I_k > 0$, задаёт изоморфизм
$\so \to \so^{*}$: вектору угловой скорости вихревой деформации
$\om = \bbI^{-1}\Om$ сопоставляется вектор завихренности
$\Om \in \so^{*} \cong \R^3$. Отождествление $\so^{*} \cong \so$
выполняется метрикой Киллинга из леммы~\ref{lem:killing}.
Случай общего положения~--- попарно различные $I_1, I_2, I_3$.
\end{definition}

Таким образом, фазовое пространство модели трёхуровневое: компактное
многообразие состояния $\RP^3 \cong \mathrm{SO}(3)$, коалгебра
$\so^{*}$ с вектором $\Om$ и инерционный изоморфизм $\om = \bbI^{-1}\Om$.
Эта структура будет использована в \S\ref{sec:ep} при выводе уравнений
движения.

\needspace{0.25\textheight}
\section{Уравнения Эйлера--Пуанкаре с диссипацией}\label{sec:ep}

\subsection{Вариационный вывод}

Зафиксируем гладкую кривую $g(t) \in \mathrm{SO}(3)$, $t \in [0,T]$,
и следом за \cite{arnold,hmr,marsdenratiu} положим
\[
  \hat{\om}(t) := g(t)^{-1}\dot g(t) \in \so,
  \qquad
  \Om(t) := \bbI\,\om(t) \in \so^{*} \cong \R^3 .
\]
Редуцированный лагранжиан~--- кинетическая энергия вихревой трубки:
$\ell(\om) = \tfrac12\langle \om, \bbI\om\rangle$; тогда
$\delta\ell/\delta\om = \Om$ (отождествление $\so^{*} \cong \so$~---
метрикой Киллинга, лемма~\ref{lem:killing}).

\begin{lemma}[об усечённых вариациях]\label{lem:variation}
Пусть $g_\varepsilon$~--- гладкая вариация кривой $g$ с закреплёнными
концами: $g_\varepsilon(0) \equiv g(0)$, $g_\varepsilon(T) \equiv g(T)$.
Положим $\hat\eta := g^{-1}\delta g$, где
$\delta g = \left.\partial_\varepsilon g_\varepsilon\right|_{\varepsilon=0}$.
Тогда $\eta(0) = \eta(T) = 0$ и
\begin{equation}\label{eq:varomega}
  \delta\hat{\om} = \dot{\hat\eta} + [\hat\om, \hat\eta] .
\end{equation}
\end{lemma}

\begin{proof}
Тождество $g_\varepsilon g_\varepsilon^{-1} \equiv \mathrm{Id}$ даёт
при дифференцировании по $\varepsilon$:
$\delta(g^{-1}) = -g^{-1}\,\delta g\,g^{-1}$. Далее, из
$\hat\om = g^{-1}\dot g$:
\[
  \delta\hat\om
  = \delta(g^{-1})\,\dot g + g^{-1}\,\delta(\dot g)
  = -g^{-1}\delta g\,g^{-1}\dot g + g^{-1}\partial_t(\delta g).
\]
Так как $\delta g = g\hat\eta$, то
$\partial_t(\delta g) = \dot g\,\hat\eta + g\,\dot{\hat\eta}$, и потому
\[
  \delta\hat\om = -\hat\eta\hat\om + \hat\om\hat\eta + \dot{\hat\eta}
  = \dot{\hat\eta} + [\hat\om,\hat\eta].
  \qedhere
\]
\end{proof}

\begin{theorem}[Эйлер--Пуанкаре]\label{th:ep}
Кривая $g(\cdot)$~--- критическая точка функционала
$\int_0^T \ell(\om)\,dt$ относительно вариаций с закреплёнными концами
тогда и только тогда, когда
\begin{equation}\label{eq:ep}
  \dot\Om = \adst_{\om}\,\Om = \Om \times \om = \Om \times \bbI^{-1}\Om .
\end{equation}
\end{theorem}

\begin{proof}
По лемме~\ref{lem:variation} и правилу дифференцирования композиции,
\[
  \delta \int_0^T \ell(\om)\,dt
  = \int_0^T \langle \Om,\, \dot\eta + \om\times\eta\rangle\,dt
  = -\int_0^T \langle \dot\Om,\,\eta\rangle\,dt
    + \int_0^T \langle \Om,\,\om\times\eta\rangle\,dt ,
\]
где граничный член $[\langle\Om,\eta\rangle]_0^T$ обнуляется, ибо
$\eta(0) = \eta(T) = 0$. Второй интеграл перепишем через коалгебраическое
действие: по определению
$\langle \adst_{\om}\Om, \eta\rangle := \langle \Om, [\om,\eta]\rangle
= \Om\cdot(\om\times\eta)$; цикличность смешанного произведения даёт
\begin{equation}\label{eq:adstar}
  \Om\cdot(\om\times\eta) = \eta\cdot(\Om\times\om)
  \qquad\Longrightarrow\qquad
  \adst_{\om}\Om = \Om\times\om .
\end{equation}
Итого $\delta S = \int_0^T \langle -\dot\Om + \Om\times\om,\;\eta\rangle dt$,
и обращение в нуль для всех $\eta$ равносильно \eqref{eq:ep}.
\end{proof}

\subsection{Диссипативное расширение и постановка задачи}

Вязкое слагаемое уравнения \eqref{eq:ns-vort} в редуцированной модели
представляется линейным оператором затухания. Далее всюду
\begin{equation}\label{eq:epd}
  \boxed{\;
  \dot\Om = \underbrace{\Om\times(\bbI^{-1}\Om)}_{\text{поворот фазы}}
  \;-\; \nu\,\bbL\,\Om,\qquad
  \Om(0) = \Om_0,}
\end{equation}
при гипотезах
\begin{itemize}[leftmargin=2.2em, itemsep=1pt]
  \item[\textup{(H1)}] $\bbL = \bbL^{\mathsf T} \succeq 0$~--- всюду;
  \item[\textup{(H2)}] $[\bbL, \bbI] = 0$ (эквивалентно: $\bbL$ диагонален
        в главном базисе $\bbI$)~--- там, где указано явно.
\end{itemize}
Слагаемое $\nu\bbL\Om$~--- структурный аналог $\nu\Delta\bwn$: оператор
$\bbL$ играет роль положительного «символа» $-\Delta$, что уточняется
балансом энстрофии в \S\ref{sec:cascade}. В покомпонентной записи
\eqref{eq:epd} принимает вид
\begin{equation}\label{eq:components}
  \begin{aligned}
    \dot\Omega_1 &= \Bigl(\tfrac1{I_3}-\tfrac1{I_2}\Bigr)\Omega_2\Omega_3 - \nu L_{11}\Omega_1,\\
    \dot\Omega_2 &= \Bigl(\tfrac1{I_1}-\tfrac1{I_3}\Bigr)\Omega_3\Omega_1 - \nu L_{22}\Omega_2,\\
    \dot\Omega_3 &= \Bigl(\tfrac1{I_2}-\tfrac1{I_1}\Bigr)\Omega_1\Omega_2 - \nu L_{33}\Omega_3,
  \end{aligned}
\end{equation}
где в последней строке использована (H2); при $\nu = 0$ это классические
уравнения Эйлера для твёрдого тела \cite{euler,arnold}.

\subsection{Гамильтонова структура: скобка Ли--Пуансона}

\begin{proposition}\label{prop:poisson}
На $C^\infty(\R^3)$ скобка
\begin{equation}\label{eq:lp}
  \{F,G\}(\Om) := -\bigl\langle \Om,\; \nabla F(\Om)\times\nabla G(\Om)\bigr\rangle
\end{equation}
корректно определена, билинейна, кососимметрична и удовлетворяет
тождеству Якоби; функция
\[
  C(\Om) = \tfrac12\|\Om\|^2
\]
является Казимиром скобки: $\{C,F\} \equiv 0$ для всех $F$. Система
\eqref{eq:ep} есть гамильтонова система $\dot f = \{f,\mathcal H\}$
с гамильтонианом $\mathcal H(\Om) = \tfrac12\langle\Om,\bbI^{-1}\Om\rangle$.
\end{proposition}

\begin{proof}
Билинейность и кососимметрия очевидны. \emph{Тождество Якоби.}
Обозначим $x = \nabla X(\Om)$, $y = \nabla Y(\Om)$, $z = \nabla Z(\Om)$,
$X'' = \nabla^2X(\Om)$ и т.д. Дифференцирование произведения в
\eqref{eq:lp} даёт (используется симметрия гессианов):
\begin{equation}\label{eq:gradbracket}
  \nabla\{X,Y\}(\Om) = -(x\times y) - X''\,(y\times\Om) - Y''\,(\Om\times x).
\end{equation}
Действительно, $\partial_i\bigl[\Om_\ell\varepsilon_{\ell mn}\partial_mX\,\partial_nY\bigr]
= \varepsilon_{imn}x_my_n + \Om_\ell\varepsilon_{\ell mn}
\bigl(X''_{im}y_n + x_mY''_{in}\bigr)$, и вторая группа сворачивается как
$\bigl[X''(y\times\Om) + Y''(\Om\times x)\bigr]_i$
(перестановка циклов смешанного произведения; симметрия $X''$, $Y''$
существенна). Подставляя \eqref{eq:gradbracket} в \eqref{eq:lp} с
внешней парой $(\,\{X,Y\}, Z\,)$ и применяя формулу Лагранжа:
\begin{equation}\label{eq:closedform}
  \{\{X,Y\},Z\}(\Om)
  = (\Om\cdot y)(x\cdot z) - (\Om\cdot x)(y\cdot z)
  + (y\times\Om)^{\mathsf T} X''\,(z\times\Om)
  + (\Om\times x)^{\mathsf T} Y''\,(z\times\Om).
\end{equation}
Просуммируем \eqref{eq:closedform} и её циклические перестановки
$(X,Y,Z) \to (Y,Z,X) \to (Z,X,Y)$. Бесгессиановая часть обнуляется:
слагаемое $(\Om\cdot v)(w\cdot u)$ входит в сумму дважды~---
с плюсом и с минусом~--- по каждому из трёх векторов. Гессиановые
члены сокращаются парами: $X''$-члены входят в
$\{\{X,Y\},Z\}$ и $\{\{Z,X\},Y\}$ как
\[
  (y\times\Om)^{\mathsf T}X''(z\times\Om)
  \;+\;(\Om\times z)^{\mathsf T}X''(y\times\Om)
  = p - p^{\mathsf T} = 0,
\]
где $p := (y\times\Om)^{\mathsf T}X''(z\times\Om)$~--- скаляр, равный
своему транспонированному виду
$p^{\mathsf T} = (z\times\Om)^{\mathsf T}X{}''^{\mathsf T}(y\times\Om)
= (z\times\Om)^{\mathsf T}X''(y\times\Om) = -(\Om\times z)^{\mathsf T}X''(y\times\Om)$,
и снова использована симметрия $X''$. Аналогично для $Y''$ и $Z''$.
Итого циклическая сумма равна нулю~--- тождество Якоби.

\emph{Казимир.} $\{C,F\} = -\langle \Om, \Om\times\nabla F\rangle
= -\det[\Om,\Om,\nabla F] = 0$.

\emph{Гамильтоновость.} Для $f \in C^\infty$ вдоль \eqref{eq:ep}:
$\dot f = \langle\nabla f, \Om\times\bbI^{-1}\Om\rangle$; с другой
стороны, $\{f,\mathcal H\} = -\langle\Om, \nabla f\times\bbI^{-1}\Om\rangle
= \langle \nabla f, \Om\times\bbI^{-1}\Om\rangle$ (цикличность).
Совпадение даёт $\dot f = \{f,\mathcal H\}$; при $\nabla f = e_k$
получаем покомпонентно \eqref{eq:ep}.
\end{proof}

\begin{remark}
Формула \eqref{eq:closedform}~--- замкнутое выражение двойной скобки
Ли--Пуансона на $\so^{*}$; стандартные доказательства (см.
\cite{marsdenratiu,arnoldkhesin}) сводят тождество Якоби к
структурным константам алгебры. Приведённое выше прямое доказательство
дополнительно даёт явный вид сокращения гессиановых членов
($p - p^{\mathsf T} = 0$).
\end{remark}

\subsection{Корректность задачи Коши}

\begin{theorem}[глобальная корректность]\label{th:wp}
При любых $\Om_0 \in \R^3$, $\nu \ge 0$ и $\bbL$, удовлетворяющей
\textup{(H1)}, задача Коши \eqref{eq:epd} имеет единственное
$C^\infty$-решение на $[0,\infty)$.
\end{theorem}

\begin{proof}
Правая часть $f(\Om) = \Om\times\bbI^{-1}\Om - \nu\bbL\Om$ гладка.
Оценка на шаре $B(0,R)$:
\[
  \|f(\Om) - f(\tilde\Om)\|
  \le \|\bbI^{-1}\|\,\bigl(\|\Om\| + \|\tilde\Om\|\bigr)\,\|\Om-\tilde\Om\|
     + \nu\|\bbL\|\,\|\Om-\tilde\Om\|
  \le \bigl(2\|\bbI^{-1}\|R + \nu\|\bbL\|\bigr)\|\Om-\tilde\Om\|,
\]
где использовано $\|a\times p - b\times q\|
\le \|(a-b)\times p\| + \|b\times(p-q)\| \le \|a-b\|\|p\| + \|b\|\|p-q\|$
и $\|\bbI^{-1}v\| \le \|\bbI^{-1}\|\,\|v\|$. По теореме
Пикара--Линделёфа существует единственное локальное решение.
Априорная оценка: вдоль решения
\[
  \frac{d}{dt}\|\Om\|^2 = 2\bigl\langle \Om,\Om\times\bbI^{-1}\Om\bigr\rangle
  - 2\nu\,\langle\Om,\bbL\Om\rangle = -2\nu\,\langle\Om,\bbL\Om\rangle \le 0
\]
(тройное произведение с двумя равными сомножителями равно нулю;
$\bbL \succeq 0$ по \textup{(H1)}). Значит $\|\Om(t)\| \le \|\Om_0\|$,
решение остаётся в компактном шаре, и критерий продолжения
\cite[гл.~II]{hartman} даёт продолжение на всё $[0,\infty)$.
\end{proof}

\needspace{0.25\textheight}
\section{Инварианты: Казимир, гамильтониан, затухание}\label{sec:inv}

\subsection{Точные законы эволюции}

\begin{theorem}[Казимир]\label{th:casimir}
Вдоль решений \eqref{eq:epd} при \textup{(H1)}
\begin{equation}\label{eq:casimir}
  \frac{d}{dt}\|\Om\|^2
  = 2\,\bigl\langle \Om,\; \Om\times\bbI^{-1}\Om\bigr\rangle
    - 2\nu\,\langle\Om,\bbL\Om\rangle
  = -2\nu\,\langle\Om,\bbL\Om\rangle \;\le\; 0 .
\end{equation}
При $\nu = 0$ величина $\|\Om(t)\|^2 \equiv \|\Om_0\|^2$ тождественно:
траектория заперта на сфере $S^2_{\|\Om_0\|} \subset \so^{*}$.
\end{theorem}

\begin{proof}
Первое равенство~--- дифференцирование квадрата нормы и подстановка
\eqref{eq:epd}. Обращение в нуль перекрёстного члена~--- определитель
с двумя равными столбцами:
$\langle\Om, \Om\times\om\rangle = \det[\Om,\Om,\om] = 0$
при любых $\om$. Знак~--- положительная полуопределённость $\bbL$
из \textup{(H1)}: $\langle\Om,\bbL\Om\rangle \ge 0$.
\end{proof}

\begin{theorem}[гамильтониан]\label{th:ham}
Вдоль решений \eqref{eq:epd}
\begin{equation}\label{eq:ham}
  \frac{d\mathcal H}{dt}
  = \bigl\langle \nabla\mathcal H,\dot\Om\bigr\rangle
  = \underbrace{\bigl\langle \bbI^{-1}\Om,\;\Om\times\bbI^{-1}\Om\bigr\rangle}_{=0}
  - \nu\,\bigl\langle \bbI^{-1}\Om,\;\bbL\Om\bigr\rangle
  = -\nu\,\Om^{\mathsf T}\bbI^{-1}\bbL\,\Om .
\end{equation}
Если дополнительно выполнена \textup{(H2)}, то
\begin{equation}\label{eq:ham-diag}
  \frac{d\mathcal H}{dt} = -\nu\sum_{k=1}^{3}\frac{L_{kk}}{I_k}\,\Omega_k^2 \;\le\; 0 ,
\end{equation}
и $\mathcal H$~--- функция Ляпунова строгого сертификата устойчивости.
\end{theorem}

\begin{proof}
Равенство \eqref{eq:ham}~--- правило цепочки с
$\nabla\mathcal H = \bbI^{-1}\Om$ и обнуление перекрёстного члена как в
теореме~\ref{th:casimir}. При \textup{(H2)} операторы $\bbI^{-1}$ и $\bbL$
диагональны в одном базисе, произведение $\bbI^{-1}\bbL$ симметрично и
положительно полуопределено, откуда \eqref{eq:ham-diag}.
\end{proof}

\begin{remark}[коммутационная тонкость]\label{rem:commutator}
Гипотеза \textup{(H2)} в теореме~\ref{th:ham} существенна и не сводится
к \textup{(H1)}: без неё квадратичная форма
$\Om \mapsto \Om^{\mathsf T}\bbI^{-1}\bbL\Om$ определяется
симметризацией $\tfrac12(\bbI^{-1}\bbL + \bbL\bbI^{-1})$, которая может
быть неопределённой даже при $\bbL \succ 0$. Пример:
\[
  \bbI^{-1} = \diag(10,1,1),\qquad
  \bbL = \begin{pmatrix} 1 & 1 & 0\\ 1 & 1{,}5 & 0\\ 0 & 0 & 1 \end{pmatrix},
  \qquad \spec \bbL = \{0{,}219;\;1;\;2{,}281\} \subset (0,\infty),
\]
\[
  \operatorname{sym}(\bbI^{-1}\bbL) = \begin{pmatrix} 10 & 5{,}5 & 0\\ 5{,}5 & 1{,}5 & 0\\ 0 & 0 & 1 \end{pmatrix},
  \qquad \det\begin{pmatrix} 10 & 5{,}5\\ 5{,}5 & 1{,}5 \end{pmatrix} = -15{,}25 < 0 .
\]
В точке $\Om_0 = (-0{,}5;\,1;\,0)$ имеем
$\Om_0^{\mathsf T}\bbI^{-1}\bbL\Om_0 = -1{,}5$, т.е.
$\dot{\mathcal H}(0) = +1{,}5\,\nu > 0$: \emph{гамильтониан возрастает,
тогда как Казимир по-прежнему строго убывает} (теорема~\ref{th:casimir}
не требует \textup{(H2)}). Численная иллюстрация~--- проверка E1 в
\S\ref{sec:numerics}: за $t \le 4$ рост $\mathcal H$ на $1{,}02\%$ при
строго убывающем $\|\Om\|^2$. Таким образом, «сертификат Ляпунова»
гарантирован только при \textup{(H2)}; энергетическая монотонность~---
свойство диагонализации, а не диссипации самой по себе.
\end{remark}

\subsection{Скорость затухания и предельное множество}

\begin{corollary}[экспоненциальное затухание]\label{cor:decay}
Если $\lambda_{\min} := \lambda_{\min}(\bbL) > 0$, то
\begin{equation}\label{eq:decay}
  \|\Om(t)\|^2 \le \|\Om_0\|^2\,e^{-2\nu\lambda_{\min} t},
  \qquad t \ge 0 .
\end{equation}
\end{corollary}

\begin{proof}
Из \eqref{eq:casimir} и неравенства Рэлея
$\langle\Om,\bbL\Om\rangle \ge \lambda_{\min}\|\Om\|^2$:
$\dot C \le -2\nu\lambda_{\min}C$ для $C = \|\Om\|^2$. Тогда
$\varphi(t) := C(t)e^{2\nu\lambda_{\min}t}$ удовлетворяет
$\dot\varphi = e^{2\nu\lambda_{\min}t}\bigl(\dot C + 2\nu\lambda_{\min}C\bigr)
\le 0$, т.е. $\varphi$ невозрастает и $\varphi(t) \le \varphi(0)$.
\end{proof}

\begin{theorem}[Ла--Салль]\label{th:lasalle}
Пусть $\omega$-предельное множество траектории \eqref{eq:epd}
непусто и компактно (это так при $\nu > 0$). Тогда оно содержится в
наибольшем инвариантном подмножестве множества
$M_0 = \{\Om : \langle\Om,\bbL\Om\rangle = 0\} = \ker\bbL$. В частности:
\begin{enumerate}[label=\textup{(\alph*)}, leftmargin=2.2em, itemsep=1pt]
  \item при $\bbL \succ 0$: $\Om(t) \to 0$ при $t\to+\infty$ для любого
        начального условия;
  \item при $\nu > 0$, диагональных $\bbL$, $\bbI$ и одномерном
        $\ker\bbL = \operatorname{span}\{e_k\}$: $\Om(t)$ сходится к
        точке равновесия $a\,e_k$, $a \in \R$.
\end{enumerate}
\end{theorem}

\begin{proof}
По теореме~\ref{th:casimir} $C = \|\Om\|^2$ невозрастает и ограничена
снизу, значит существует $C_\infty = \lim C(t)$. На $\omega$-предельном
множестве $C \equiv C_\infty$ и $\dot C = 0$; по принципу
Ла--Салля \cite{lasalle} $\omega$-предельное множество инвариантно и
лежит в $M_0 = \ker\bbL$. (a): $\ker\bbL = \{0\}$, и единственная
инвариантная точка в $\{0\}$~--- само $0$ (равновесие). (b): орбита,
целиком лежащая в $\operatorname{span}\{e_k\}$, имеет вид
$\Om(t) = a(t)e_k$; подстановка в \eqref{eq:epd} даёт
$\dot a\,e_k = a^2\,e_k\times(e_k/I_k) = 0$, откуда $a$ постоянна:
инвариантные точки в $\ker\bbL$~--- в точности равновесия $ae_k$.
\end{proof}

\subsection{Топология орбит при \texorpdfstring{$\nu=0$}{nu=0}}

\begin{proposition}[периодичность]\label{prop:periodic}
Пусть $\nu = 0$, моменты $I_1, I_2, I_3$ попарно различны и
$\Om_0$ не лежит на осях $\{ \pm a e_k \}$. Тогда орбита
$\{\Om(t)\}$~--- замкнутая гладкая кривая (окружность Лиувилля),
и движение по ней периодично.
\end{proposition}

\begin{proof}
\emph{Шаг 1: равновесия.} При $\nu = 0$ равновесия удовлетворяют
$\Om\times\bbI^{-1}\Om = 0$, т.е. $\Om \parallel \bbI^{-1}\Om$~---
$\Om$ собственный вектор $\bbI^{-1}$, значит $\Om \in \bigcup_k
\operatorname{span}\{e_k\}$; на сфере $S^2_R$ это шесть точек
$\pm Re_k$, $R = \|\Om_0\|$.

\emph{Шаг 2: критические точки $\mathcal H|_{S^2_R}$.} Множитель
Лагранжа для условного экстремума $\mathcal H$ на сфере даёт
$\Omega_k(1/I_k - \lambda) = 0$ для всех $k$; нетривиальные решения~---
снова $\pm Re_k$. В этих шести точках $\mathcal H|_{S^2_R}$ принимает
три попарно различных значения $\tfrac{R^2}{2I_k}$, и все они
невырождены (попарное различие $I_k$).

\emph{Шаг 3: структура регулярных уровней.} Пусть
$h \notin \{\tfrac{R^2}{2I_k}\}$ и
$L_h = \{\Om \in S^2_R : \mathcal H(\Om) = h\}$. На $L_h$ градиент
$\mathcal H|_{S^2_R}$ не обращается в нуль, поэтому по теореме о
неявной функции $L_h$~--- гладкое компактное одномерное многообразие
без края; по классификации компактных одномерных многообразий каждая
его компонента диффеоморфна окружности. По теореме~\ref{th:casimir} и
предложению~\ref{prop:poisson} ($\nu = 0$: гамильтонов поток) обе
функции $C$ и $\mathcal H$ сохраняются, значит каждая компонента
$L_h$ инвариантна относительно потока.

\emph{Шаг 4: периодичность.} На $L_h$ векторное поле \eqref{eq:ep} не
обращается в нуль (нули~--- в точности равновесия шага~1, лежащие на
критических уровнях). Параметризуем компоненту углом $\theta$; тогда
$\dot\theta = a(\theta)$, где $a$ непрерывна и без нулей на
компактной окружности, значит $|a| \ge a_0 > 0$ и время полного
обхода
\[
  T = \oint \frac{d\theta}{|a(\theta)|} \le \frac{2\pi}{a_0} < \infty .
\]
Следовательно каждая орбита~--- вся компонента $L_h$, проходимая за
конечное время $T$: движение периодично.
\end{proof}

\begin{proposition}[устойчивость равновесий]\label{prop:stability}
Пусть $\nu = 0$, $I_1 > I_2 > I_3$. На сфере $S^2_R$ точки $\pm Re_1$
и $\pm Re_3$ устойчивы по Ляпунову, точки $\pm Re_2$ неустойчивы:
линеаризация в $Re_2$ имеет действительные собственные значения
\[
  \lambda_{\pm} = \pm R\sqrt{\Bigl(\tfrac1{I_1}-\tfrac1{I_2}\Bigr)
                           \Bigl(\tfrac1{I_2}-\tfrac1{I_3}\Bigr)} \ne 0 .
\]
\end{proposition}

\begin{proof}
\emph{Устойчивость $\pm Re_1$, $\pm Re_3$.} На $S^2_R$ вблизи $Re_1$
($I_1$~--- максимум)
\[
  \mathcal H - \mathcal H(Re_1)
  = \tfrac12\sum_{k=2}^{3}\Bigl(\tfrac1{I_k}-\tfrac1{I_1}\Bigr)\Omega_k^2
  \ge c\,\bigl(\Omega_2^2+\Omega_3^2\bigr), \qquad c > 0,
\]
т.е. $\mathcal H|_{S^2_R}$ имеет строгий локальный минимум в $Re_1$;
сохранение $\mathcal H$ и тот факт, что малые уровни $\mathcal H|_{S^2_R}$
суть малые замкнутые кривые вокруг $Re_1$ (следствие
предложения~\ref{prop:periodic}), дают устойчивость по Ляпунову;
для $Re_3$ ($I_3$~--- минимум) аналогично с максимумом $-\mathcal H$.

\emph{Неустойчивость $Re_2$.} Положим $\Om = Re_2 + \xi$,
$\|\xi\|$ мало. Линеаризация \eqref{eq:ep}:
\[
  \dot\xi = Re_2\times\bbI^{-1}\xi + \xi\times\bbI^{-1}(Re_2)
  = R\Bigl(\tfrac{\xi_3}{I_3},\,0,\,-\tfrac{\xi_1}{I_1}\Bigr)
    + \Bigl(-\tfrac{R}{I_2}\xi_3,\,0,\,\tfrac{R}{I_2}\xi_1\Bigr),
\]
и после приведения
\[
  \dot\xi_1 = -R\Bigl(\tfrac1{I_2}-\tfrac1{I_3}\Bigr)\xi_3,
  \qquad
  \dot\xi_3 = R\Bigl(\tfrac1{I_2}-\tfrac1{I_1}\Bigr)\xi_1,
  \qquad \dot\xi_2 = 0 .
\]
Матрица $\begin{psmallmatrix}0 & -a\\ b & 0\end{psmallmatrix}$,
$a = R(\tfrac1{I_2}-\tfrac1{I_3}) < 0$,
$b = R(\tfrac1{I_2}-\tfrac1{I_1}) > 0$, имеет
$\lambda_{\pm}^2 = -ab > 0$, т.е. $\lambda_+ > 0$~--- по теореме об
неустойчивости по первому приближению \cite[гл.~VII]{hartman}
равновесие неустойчиво (эффект «теннисной ракетки»/Джанибекова).
\end{proof}

\needspace{0.25\textheight}
\section{Каскад редукции \texorpdfstring{$P^3 \to 2D \to 1D \to 0D$}{P3 -> 2D -> 1D -> 0D}}\label{sec:cascade}

\subsection{Шаг $P^3 \to 2D$: обнуление растяжения}

\begin{lemma}[о растяжении]\label{lem:strain}
Пусть $\bu \in C^1$, $\nabla\cdot\bu = 0$, $\bwn = \nabla\times\bu$. Обозначим
\[
  S := \tfrac12\bigl(\nabla\bu + (\nabla\bu)^{\mathsf T}\bigr),\qquad
  W := \tfrac12\bigl(\nabla\bu - (\nabla\bu)^{\mathsf T}\bigr)
\]
тензоры деформации и завихривания. Тогда
\begin{equation}\label{eq:stretch}
  (\bwn\cdot\nabla)\bu = S\,\bwn
  \qquad\text{и}\qquad
  W\,\bwn = 0 .
\end{equation}
\end{lemma}

\begin{proof}
В компонентах $W_{ij} = \tfrac12(\partial_ju_i - \partial_iu_j)
= \tfrac12\varepsilon_{ijk}\,\omega_k$ (проверяется на паре $(i,j)=(1,2)$
и распространяется циклическими перестановками). Тогда
\[
  (W\bwn)_i = W_{ij}\omega_j = \tfrac12\varepsilon_{ijk}\omega_k\omega_j = 0 ,
\]
поскольку сумма $\varepsilon_{ijk}\omega_k\omega_j$ переходит в себя
при переименовании $j \leftrightarrow k$ и одновременно меняет знак.
Далее, $(\bwn\cdot\nabla)\bu = (\nabla\bu)\,\bwn = (S+W)\bwn = S\bwn$.
\end{proof}

\begin{theorem}[Ладыженская, энстрофия]\label{th:ladyzh}
Пусть $\bu = \bu(x_1,x_2,t)$ двуммерно и бездивергентно. Тогда
$\bwn = \omega_3\,e_3$, член растяжения тождественно равен нулю:
$(\bwn\cdot\nabla)\bu = \omega_3\,\partial_3\bu = 0$, и энстрофия
удовлетворяет
\begin{equation}\label{eq:enstrophy}
  \frac{d}{dt}\,\frac12\|\bwn(\cdot,t)\|_{L^2}^2 = -\nu\,\|\nabla\bwn(\cdot,t)\|_{L^2}^2 \;\le\; 0 .
\end{equation}
\end{theorem}

\begin{proof}
Умножим \eqref{eq:ns-vort} скалярно на $\bwn$ и проинтегрируем по
периодической области (или по $\R^2$ с убыванием на бесконечности):
\[
  \tfrac12\frac{d}{dt}\|\bwn\|_{L^2}^2
  = \int (\bwn\cdot\nabla)\bu\cdot\bwn\,dx
  - \int(\bu\cdot\nabla)\bwn\cdot\bwn\,dx
  + \nu\int \Delta\bwn\cdot\bwn\,dx .
\]
Транспортный член обнуляется интегрированием по частям
($\nabla\cdot\bu = 0$): $\int(\bu\cdot\nabla)\bwn\cdot\bwn\,dx
= -\tfrac12\int(\bu\cdot\nabla)|\bwn|^2dx = 0$. Член растяжения равен
нулю по условию теоремы. Вязкий член:
$\int\Delta\bwn\cdot\bwn\,dx = -\|\nabla\bwn\|_{L^2}^2$.
\end{proof}

\begin{proposition}[2D-сечение модели]\label{prop:2dplane}
Плоскость $\{\Omega_3 = 0\}$ инвариантна относительно \eqref{eq:epd}
тогда и только тогда, когда $I_1 = I_2$. В этом случае на сечении
\[
  \dot\Omega_1 = -\nu L_{11}\Omega_1,\qquad
  \dot\Omega_2 = -\nu L_{22}\Omega_2,\qquad
  \frac{d}{dt}\bigl(\Omega_1^2+\Omega_2^2\bigr)
  = -2\nu\bigl(L_{11}\Omega_1^2 + L_{22}\Omega_2^2\bigr) \le 0,
\]
что в точности повторяет структуру \eqref{eq:enstrophy}.
\end{proposition}

\begin{proof}
По \eqref{eq:components} на плоскости $\Omega_3 = 0$
\[
  \dot\Omega_3 = \Bigl(\tfrac1{I_2}-\tfrac1{I_1}\Bigr)\Omega_1\Omega_2 ,
\]
и это выражение тождественно нулю на всей плоскости (при произвольных
$\Omega_1, \Omega_2$) iff $I_1 = I_2$; необходимость очевидна
($\Omega_1\Omega_2 \not\equiv 0$), достаточность~--- подстановка.
Закон для $C_2 = \Omega_1^2+\Omega_2^2$~--- прямое вычисление.
\end{proof}

\subsection{Шаг $2D \to 1D$: линеаризация Коула--Хопфа}

\begin{theorem}[Коул--Хопф]\label{th:colehopf}
Пусть $\phi \in C^{2,1}$, $\phi > 0$, и $\phi$ решает уравнение
теплопроводности $\partial_t\phi = \nu\,\partial_{xx}\phi$ на окружности
$\mathbb T$. Тогда
\begin{equation}\label{eq:colehopf}
  u(x,t) := -2\nu\,\partial_x\ln\phi(x,t) = -2\nu\,\frac{\phi_x}{\phi}
\end{equation}
решает уравнение Бюргерса $u_t + uu_x = \nu u_{xx}$ на $\mathbb T$.
\end{theorem}

\begin{proof}
Обозначим $B := (\ln\phi)_x$, $A := (\ln\phi)_{xx} = B_x$, так что
$u = -2\nu B$, $u_x = -2\nu A$, $u_{xx} = -2\nu A_x$,
$u_t = -2\nu B_t$. Из $\phi_t = \nu\phi_{xx}$:
\[
  B_t = \partial_x\Bigl(\frac{\phi_t}{\phi}\Bigr)
      = \partial_x\Bigl(\frac{\nu\phi_{xx}}{\phi}\Bigr)
      = \nu\,\partial_x\bigl(A + B^2\bigr)
      = \nu A_x + 2\nu B B_x = \nu A_x + 2\nu B A .
\]
Подставим в невязку уравнения Бюргерса:
\[
  u_t + uu_x - \nu u_{xx}
  = -2\nu B_t + 4\nu^2 BA + 2\nu^2 A_x
  = -2\nu\bigl(\nu A_x + 2\nu BA\bigr) + 4\nu^2 BA + 2\nu^2 A_x = 0.
  \qedhere
\]
\end{proof}

\begin{corollary}[глобальная регулярность]\label{cor:burgers-global}
Для гладкого $\phi_0 > 0$ на $\mathbb T$ существует единственное
гладкое решение Бюргерса при всех $t > 0$: по принципу максимума
$\phi(\cdot,t) \ge \min \phi_0 > 0$, тёплое ядро даёт
\[
  \phi(x,t) = \int_{\mathbb T} G(x-y,t)\,\phi_0(y)\,dy \in C^\infty,
  \qquad
  \|u(\cdot,t)\|_\infty \le \frac{2\nu\,\|\phi_x(\cdot,t)\|_\infty}{\min_{\mathbb T}\phi_0},
\]
и сингулярности (образование ударной волны при $\nu = 0$) полностью
регуляризуются вязкостью.
\end{corollary}

\begin{proof}
Единственность и гладкость $\phi$~--- стандартная теория уравнения
теплопроводности на компактном многообразии; положительность~---
принцип максимума; формула~--- свёртка с ядром Грина; оценка~---
непосредственно из \eqref{eq:colehopf}.
\end{proof}

\subsection{Шаг $1D \to 0D$: спектральный проектор МакВини}

\begin{definition}\label{def:proj}
Для $\Om \ne 0$ положим
\[
  \bbP(\Om) := \frac{\Om\,\Om^{\mathsf T}}{\|\Om\|^2} \in \R^{3\times3},
  \qquad
  \mathcal Q(\bbP) := 3\bbP^2 - 2\bbP^3 .
\]
\end{definition}

\begin{lemma}[точный проектор]\label{lem:rank1}
При любом $\Om \ne 0$ матрица $\bbP(\Om)$~--- ортогональный проектор:
$\bbP^2 = \bbP = \bbP^{\mathsf T}$, $\spec\bbP = \{1,0,0\}$, и
$\mathcal Q(\bbP) = \bbP$ тождественно.
\end{lemma}

\begin{proof}
\[
  \bbP^2 = \frac{\Om(\Om^{\mathsf T}\Om)\Om^{\mathsf T}}{\|\Om\|^4}
  = \frac{\Om\|\Om\|^2\Om^{\mathsf T}}{\|\Om\|^4} = \bbP;
  \qquad
  \bbP\,\Om = \frac{\Om(\Om^{\mathsf T}\Om)}{\|\Om\|^2} = \Om,
\]
а на ортогональном дополнении $\operatorname{span}\{\Om\}^\perp$ имеем
$\bbP v = 0$; отсюда спектр. Подстановка $\bbP^2 = \bbP$,
$\bbP^3 = \bbP$ даёт $\mathcal Q(\bbP) = 3\bbP - 2\bbP = \bbP$.
\end{proof}

\begin{lemma}[спектральный расчёт]\label{lem:q-spectral}
Пусть $\bbP = \bbP^{\mathsf T}$, $\spec\bbP \subset [0,1]$, и
$q(\lambda) = 3\lambda^2 - 2\lambda^3$. Тогда
\begin{equation}\label{eq:factorization}
  q(\lambda)^2 - q(\lambda)
  = -\lambda^2(1-\lambda)^2\,(3-2\lambda)(2\lambda+1),
\end{equation}
и, следовательно,
\begin{equation}\label{eq:q-bound}
  \|\mathcal Q(\bbP)^2 - \mathcal Q(\bbP)\|
  \le 4\,\|\bbP^2 - \bbP\|^2 .
\end{equation}
\end{lemma}

\begin{proof}
Разложение \eqref{eq:factorization}: $q(\lambda) = \lambda^2(3-2\lambda)$;
далее $(1-\lambda)^2(2\lambda+1) = 1 - 3\lambda^2 + 2\lambda^3$, откуда
$q(\lambda) - 1 = -(1-\lambda)^2(2\lambda+1)$, и перемножение даёт
\eqref{eq:factorization}. На отрезке $[0,1]$ функция
$p(\lambda) := (3-2\lambda)(2\lambda+1) = -4\lambda^2+4\lambda+3$
удовлетворяет $p(\lambda) \le 4$, причём $p(\tfrac12) = 4$~--- максимум
($p'(\lambda) = -8\lambda+4$). Поэтому
\[
  |q(\lambda)^2 - q(\lambda)| = \lambda^2(1-\lambda)^2\,p(\lambda)
  \le 4\lambda^2(1-\lambda)^2 .
\]
Функциональное исчисление: матрица $\bbP$ эрмитова, значит унитарно
диагонализуема, $\bbP = U D U^{*}$, где $D = \diag(\lambda_1,\lambda_2,\lambda_3)$;
полиномы от $\bbP$ диагонализуются тем же $U$, и
\begin{gather*}
  \|\bbP^2-\bbP\| = \max_i|\lambda_i^2-\lambda_i| = \max_i \lambda_i(1-\lambda_i),\\
  \|\mathcal Q(\bbP)^2-\mathcal Q(\bbP)\|
  = \max_i |q(\lambda_i)^2-q(\lambda_i)| ,
\end{gather*}
так как $\mathcal Q(\bbP)^2-\mathcal Q(\bbP) = (q^2-q)(\bbP)$ также
эрмитова. Комбинируя с поточечной оценкой:
\[
  \|\mathcal Q(\bbP)^2-\mathcal Q(\bbP)\|
  \le 4\max_i \lambda_i^2(1-\lambda_i)^2
  \le 4\Bigl(\max_i \lambda_i(1-\lambda_i)\Bigr)^{\!2}
  = 4\,\|\bbP^2-\bbP\|^2 . \qedhere
\]
\end{proof}

\begin{theorem}[оценка дефекта МакВини; оптимальность константы]\label{th:mcweeny}
Оценка \eqref{eq:q-bound} верна для всех эрмитовых $\bbP$ со спектром
в $[0,1]$, и константа $4$ неулучшаема: для скалярного случая
$\bbP = (\tfrac12)$ имеет место точное равенство
\[
  \|\mathcal Q(\bbP)^2-\mathcal Q(\bbP)\| = \tfrac14 = 4\cdot\Bigl(\tfrac14\Bigr)^{\!2} = 4\,\|\bbP^2-\bbP\|^2 .
\]
\end{theorem}

\begin{proof}
Оценка~--- лемма~\ref{lem:q-spectral}. Равенство при $\lambda=\tfrac12$:
$\lambda(1-\lambda) = \tfrac14$,
$|q^2-q| = \tfrac14\cdot\tfrac14\cdot p(\tfrac12) = \tfrac1{16}\cdot4 = \tfrac14$
и $4\|\bbP^2-\bbP\|^2 = 4\cdot\tfrac1{16} = \tfrac14$.
\end{proof}

\begin{proposition}[итерация МакВини]\label{prop:mcweeny-iter}
Точки $\lambda \in \{0,\,1\}$~--- притягивающие неподвижные точки
отображения $q$, точка $\tfrac12$~--- отталкивающая:
$q'(\tfrac12) = \tfrac32 > 1$. Для $\lambda_0 \in [0,1)$ итерации
$\lambda_{n+1} = q(\lambda_n)$ монотонны и сходятся к $0$ при
$\lambda_0 \in (0,\tfrac12)$ и к $1$ при $\lambda_0 \in (\tfrac12,1)$.
\end{proposition}

\begin{proof}
Разложение
\[
  q(\lambda) - \lambda = \lambda(1-\lambda)(2\lambda-1)
\]
проверяется раскрытием скобок и даёт знак: $q(\lambda) < \lambda$ на
$(0,\tfrac12)$, $q(\lambda) > \lambda$ на $(\tfrac12,1)$. Монотонная
ограниченная последовательность сходится; предел~--- неподвижная точка
$q$, а неподвижные точки суть корни $\lambda(1-\lambda)(2\lambda-1)=0$,
т.е. $\{0,\tfrac12,1\}$; предел не может равняться $\tfrac12$ для
строго монотонной последовательности. Производная
$q'(\lambda) = 6\lambda-6\lambda^2$, $q'(\tfrac12) = \tfrac32$.
\end{proof}

\begin{remark}
В самой модели $\bbP(\Om(t))$~--- \emph{точно} идемпотентный проектор
(лемма~\ref{lem:rank1}) при всех $t$ с $\Om(t) \ne 0$: дефект
идемпотентности равен нулю тождественно, и $0 \le 4\cdot 0^2$~--- оценка
\eqref{eq:q-bound} вырождается в равенство нулей. Содержательная роль
полинома МакВини~--- устойчивость к возмущениям: если численная или
модельная погрешность выводит матрицу из множества проекторов
($\|\bbP^2-\bbP\| = \varepsilon$), однократное применение $\mathcal Q$
подавляет дефект квадратично: до величины $\le 4\varepsilon^2$
(см.~\cite{mcweeny}). Особенность $\Om = 0$ устранима: $\bbP$ не
определена лишь в точке, куда траектория приходит только
асимптотически при $\bbL \succ 0$ (теорема~\ref{th:lasalle}).
\end{remark}

\needspace{0.25\textheight}
\section{Отсутствие сингулярностей в модели}\label{sec:noblowup}

\begin{theorem}[главная]\label{th:main}
Пусть выполнена \textup{(H1)}. Тогда для решения \eqref{eq:epd}:
\begin{enumerate}[label=\textup{(\roman*)}, leftmargin=2.4em, itemsep=1pt]
  \item решение существует и единственно при всех $t \ge 0$ при любом
        $\nu \ge 0$ (теорема~\ref{th:wp});
  \item \emph{априорная оценка}: $\displaystyle\sup_{t\ge0}\|\Om(t)\|
        \le \|\Om_0\| < \infty$ (теорема~\ref{th:casimir});
  \item при $\nu = 0$: $\|\Om(t)\| \equiv \|\Om_0\|$, орбита замкнута
        и периодична, если $\Om_0$ не осевая
        (предложение~\ref{prop:periodic});
  \item при $\nu > 0$, $\lambda_{\min}(\bbL) > 0$:
        $\|\Om(t)\| \le \|\Om_0\|\,e^{-\nu\lambda_{\min}t} \to 0$
        (следствие~\ref{cor:decay}, теорема~\ref{th:lasalle});
  \item при \textup{(H2)}: $\mathcal H(t) \le \mathcal H(0)$
        (теорема~\ref{th:ham});
  \item \emph{нет сингулярности}: для любого $T^{*} \le \infty$
        \begin{equation}\label{eq:noblowup}
          \limsup_{t \to T^{*}} \|\Om(t)\| \;\le\; \|\Om_0\| \;<\; \infty ,
        \end{equation}
        и все «серриновские» интегралы конечны тривиально:
        $\int_0^{T}\|\Om(t)\|^{p}dt \le \|\Om_0\|^{p}T < \infty$.
\end{enumerate}
\end{theorem}

\begin{proof}
Пункты (i)--(v)~--- сборка результатов \S\ref{sec:ep}--\S\ref{sec:inv};
покажем (vi). По (ii) функция $t \mapsto \|\Om(t)\|$ ограничена
константой $\|\Om_0\|$ на всём $[0,\infty)$, откуда
$\limsup_{t\to T^{*}}\|\Om(t)\| \le \|\Om_0\|$ для любого
$T^{*} \le \infty$. Конечность интегралов~--- непосредственная
оценка под интегралом.
\end{proof}

\begin{proof}[Комментарий к доказательству]
Содержание оценки \eqref{eq:noblowup}~--- в её механизме: вся
нелинейность $\Om\times\bbI^{-1}\Om$ есть чисто кососимметричный
поворот фазы, ортогональный вектору состояния
($\langle\Om,\Om\times\om\rangle = \det[\Om,\Om,\om] = 0$ при всех
$\om$), а потому \emph{в принципе не способна увеличивать норму}
завихренности; всё возрастание подавлено структурой, а не малостью
коэффициентов. Знак определяет только диссипативный член, и он
неотрицателен по \textup{(H1)}.
\end{proof}

\subsection{Явные решения: осесимметричный случай}

\begin{proposition}[явные решения]\label{prop:axisym}
Пусть $I_1 = I_2 =: I_\perp$ и $\bbL = \diag(\lambda,\lambda,L_{33})$
изотропен в плоскости симметрии. Положим
$z(t) := \Omega_1(t) + i\,\Omega_2(t)$,
$z_0 := \Omega_1(0) + i\,\Omega_2(0)$,
$\Omega_3^0 := \Omega_3(0)$,
$\kappa_0 := \bigl(\tfrac1{I_3}-\tfrac1{I_\perp}\bigr)\Omega_3^0$.
Тогда решение \eqref{eq:epd} выражается в явном виде:
\begin{align}
  \nu = 0:\qquad & \Omega_3(t) \equiv \Omega_3^0, \qquad
    z(t) = z_0\,e^{-i\kappa_0 t}, \label{eq:exact-nu0}\\
  \nu > 0,\ L_{33} > 0:\qquad & \Omega_3(t) = \Omega_3^0\,e^{-\nu L_{33}t},
    \qquad z(t) = z_0\,
    \exp\!\Bigl(-\nu\lambda t - i\,\kappa_0\,
    \frac{1-e^{-\nu L_{33}t}}{\nu L_{33}}\Bigr). \label{eq:exact-nup}
\end{align}
\end{proposition}

\begin{proof}
При $I_1 = I_2$ система \eqref{eq:components} принимает вид
\[
  \dot\Omega_1 = \kappa(t)\,\Omega_2 - \nu\lambda\,\Omega_1,\qquad
  \dot\Omega_2 = -\kappa(t)\,\Omega_1 - \nu\lambda\,\Omega_2,\qquad
  \dot\Omega_3 = -\nu L_{33}\,\Omega_3,
\]
где $\kappa(t) = \bigl(\tfrac1{I_3}-\tfrac1{I_\perp}\bigr)\Omega_3(t)$
(перекрёстные члены $\propto \Omega_1\Omega_3$, $\Omega_2\Omega_3$
симметричны и сокращаются). Третье уравнение~--- линейное, откуда
$\Omega_3(t) = \Omega_3^0e^{-\nu L_{33}t}$ (при $\nu = 0$~---
константа). Для $z$: $\dot z = \dot\Omega_1 + i\dot\Omega_2
= \kappa\Omega_2 - \nu\lambda\Omega_1 + i(-\kappa\Omega_1 - \nu\lambda\Omega_2)
= -i\kappa(t)\,z - \nu\lambda\,z$. Это линейное уравнение с переменным
коэффициентом, и
\[
  z(t) = z_0\exp\!\Bigl(\int_0^t\bigl(-i\kappa(s) - \nu\lambda\bigr)ds\Bigr),
  \qquad
  \int_0^t \kappa(s)\,ds = \kappa_0\,\frac{1-e^{-\nu L_{33}t}}{\nu L_{33}},
\]
что даёт \eqref{eq:exact-nup}; при $\nu = 0$ интеграл равен $\kappa_0 t$,
и получается \eqref{eq:exact-nu0}.
\end{proof}

\begin{remark}[конечная фаза]
Формула \eqref{eq:exact-nup} вскрывает тонкий эффект: суммарный угол
поворота фазы за бесконечное время конечен,
\[
  \int_0^\infty \kappa(s)\,ds = \frac{\kappa_0}{\nu L_{33}} < \infty :
\]
вязкое затухание компоненты $\Omega_3$ останавливает прецессию.
При $\nu = 0$ движение~--- чистая прецессия с периодом
$T = 2\pi/|\kappa_0|$ (частный случай
предложения~\ref{prop:periodic}). Численная проверка точности формул
\eqref{eq:exact-nu0}--\eqref{eq:exact-nup} и порядка сходимости RK4~---
тесты A1--A2 в \S\ref{sec:numerics}.
\end{remark}

\needspace{0.25\textheight}
\section{Критический разбор: границы переноса на полное уравнение}\label{sec:critique}

\subsection{Структурный словарь}

Сведём соответствие между моделью \eqref{eq:epd} и полным уравнением
\eqref{eq:ns-vort} (см.~\cite{arnoldkhesin,marsdenratiu}):

\begin{table}[htbp]
\centering
\small
\caption{Структурный словарь «модель $\leftrightarrow$ полное уравнение».}
\label{tab:dictionary}
\begin{tabularx}{\textwidth}{@{}l X l@{}}
\toprule
Модель & Содержание & Полное уравнение \\
\midrule
$\Om \in \so^{*}$ & амплитудный вектор вихревой трубки &
  $\bwn(x)$, бесконечномерно \\
$\Om\times\bbI^{-1}\Om$ & коалгебраический транспорт (коммутатор) &
  $-[u,\bwn] = \nabla\times(u\times\bwn)$ \\
$-\nu\bbL\Om$, $\bbL\succeq0$ & диссипация через положительный символ &
  $\nu\Delta\bwn$ \\
$C = \|\Om\|^2$ & Казимир & $\tfrac12\|\bwn\|_{L^2}^2$ (энстрофия) \\
$\mathcal H$ & энергия & $\tfrac12\|u\|_{L^2}^2$ \\
$S^2_{\|\Om_0\|} \subset \so^{*}$ & коалгебраическая орбита &
  орбиты группы диффеоморфизмов \\
\bottomrule
\end{tabularx}
\end{table}

\subsection{Почему модель работает: сферичность орбит}

\begin{proposition}[орбиты $\so^{*}$]\label{prop:orbits}
Коалгебраические орбиты действия $\mathrm{SO}(3)$ на $\so^{*}
\cong \R^3$~--- в точности сферы $\{\|\Om\| = \mathrm{const}\}$.
Следовательно любая радиальная функция, в частности $C = \|\Om\|^2$,~---
Казимир скобки \eqref{eq:lp}, и сохранение $\|\Om\|$ при $\nu=0$
является \emph{чисто геометрическим} фактом.
\end{proposition}

\begin{proof}
Инфинитезимальное коалгебраическое действие вычислено в
\eqref{eq:adstar}: $\adst_{\xi}\Om = \Om\times\xi$. Касательное
пространство орбиты в точке $\Om$ есть
$\{\Om\times\xi : \xi \in \R^3\} = (\operatorname{span}\{\Om\})^{\perp}$~--- касательная плоскость к сфере радиуса $\|\Om\|$ в точке $\Om$.
Значит орбита~--- открытое подмногообразие соответствующей сферы; будучи
орбитой действия компактной группы, она компактна и потому совпадает
со всей сферой (сфера связна). Радиальные функции постоянны на орбитах,
т.е. коммутируют со скобкой \eqref{eq:lp} в смысле
предложения~\ref{prop:poisson}.
\end{proof}

В этом~--- точное структурное объяснение всей конструкции: сферичность
орбит \emph{гарантирует} закон Казимира. Для двумерной идеальной
жидкости аналогичное утверждение истинно: энстрофия
$\tfrac12\|\bwn\|_{L^2}^2$~--- подлинный Казимир бесконечномерной
скобки Ли--Пуансона группы сохраняющих площадь диффеоморфизмов, и
теорема Ладыженской \eqref{eq:enstrophy} есть её вязкий аналог
\cite{arnoldkhesin,ladyzh}. Для трёхмерной идеальной жидкости Казимиром
служит только гелисичность, а не энстрофия: орбиты несферичны, и член
растяжения $(\bwn\cdot\nabla)\bu = S\bwn$ (лемма~\ref{lem:strain})
выводит норму $\|\bwn\|_{L^2}$ с орбитально-инвариантных множеств.
Баланс энстрофии полного уравнения
\begin{equation}\label{eq:full-enstrophy}
  \frac{d}{dt}\,\frac12\|\bwn\|_{L^2}^2
  = \int \langle S\bwn, \bwn\rangle\,dx \;-\; \nu\|\nabla\bwn\|_{L^2}^2
\end{equation}
содержит штрейн-член $\int\langle S\bwn,\bwn\rangle dx$, который модель
обнуляет \emph{по построению}: в \eqref{eq:epd} ему просто негде жить~---
кососимметричная нелинейность $\Om\times\bbI^{-1}\Om$ структурно
не способна его породить (теорема~\ref{th:main}).

\subsection{Условная теорема о тензоре деформации}

Сделаем границу метода количественной. Расширим модель явным
штрейн-членом:
\begin{equation}\label{eq:epd-strain}
  \dot\Om = \Om\times\bbI^{-1}\Om + S(t)\,\Om - \nu\,\bbL\,\Om,
\end{equation}
где $S(t) = S(t)^{\mathsf T}$~--- произвольные ограниченные измеримые
симметричные матрицы (экзогенный тензор деформации).

\begin{theorem}[условная]\label{th:strain}
Вдоль решений \eqref{eq:epd-strain} при \textup{(H1)} и
$\lambda_{\min} := \lambda_{\min}(\bbL) > 0$
\begin{equation}\label{eq:strain-bound}
  \|\Om(t)\| \;\le\; \|\Om_0\|\,
  \exp\!\Bigl(\int_0^t \|S(s)\|\,ds \;-\; \nu\lambda_{\min}\,t\Bigr)
  \qquad \forall\, t \ge 0 .
\end{equation}
В частности, если штрейн интегрируем с запасом:
$\int_0^t\|S(s)\|ds \le \nu\lambda_{\min}t + K$ на $[0,T]$, то
$\|\Om(t)\| \le e^{K}\|\Om_0\|$~--- сингулярность невозможна.
\end{theorem}

\begin{proof}
Дифференцируем $C = \|\Om\|^2$ вдоль \eqref{eq:epd-strain}:
\[
  \dot C = 2\underbrace{\langle\Om,\Om\times\bbI^{-1}\Om\rangle}_{=0}
  + 2\langle S(t)\Om,\Om\rangle - 2\nu\langle\Om,\bbL\Om\rangle
  \le \bigl(2\|S(t)\| - 2\nu\lambda_{\min}\bigr)\,C ,
\]
где обнуление перекрёстного члена~--- как в теореме~\ref{th:casimir},
а неравенство~--- Коши--Буняковский с $\|S\| := \|S\|_{2}$ и оценка
Рэлея для $\bbL$. Дифференциальное неравенство
$\dot C \le \beta(t)C$ интегрируется (Грёнуолл):
$C(t) \le C(0)\exp\int_0^t\beta(s)ds$, что после извлечения корня даёт
\eqref{eq:strain-bound}. Следствие о $e^K\|\Om_0\|$~--- подстановка
условия интегрируемости.
\end{proof}

\begin{corollary}[аналогия с BKM/CKN]\label{cor:bkm}
Структура теоремы~\ref{th:strain}~--- точный аналог классических
критериев регулярности полного уравнения: критерий
Бил--Като--Маджды связывает отсутствие сингулярности на $[0,T]$ с
конечностью $\int_0^T\|\bwn\|_{L^\infty}dt$ \cite{bkm}, критерий
Каффарелли--Кона--Ниренберга~--- с малостью интеграла от
характеристики $A(t) \sim \bigl(1 + \|\nabla\bu\|_{L^\infty}\bigr)$
\cite{ckn}. В модели роль $\|\bwn\|_{L^\infty}$ играет $\|S(t)\|$:
\emph{регулярность определяется интегрируемой силой растяжения}.
Численно оценка \eqref{eq:strain-bound} точна: при изотропном
$S \equiv s_0\,\mathrm{Id}$ выполняется $\|\Om(t)\| =
\|\Om_0\|e^{(s_0-\nu\lambda_{\min})t}$ с \emph{равенством} в
\eqref{eq:strain-bound} (тест S1b в \S\ref{sec:numerics}), а при
$s_0 > \nu\lambda_{\min}$ реализуется экспоненциальный рост~--- модель
с несжимаемым штрейном не имеет априорной оценки, что и есть
честная граница метода.
\end{corollary}

\subsection{Что доказано и что остаётся открытым}

\begin{enumerate}[label=\textup{(\arabic*)}, leftmargin=2.2em, itemsep=2pt]
  \item \emph{Доказано строго} (внутри модели): все результаты
        \S\ref{sec:geom}--\S\ref{sec:noblowup}~--- геометрия
        $\RP^3\cong\mathrm{SO}(3)$, вывод \eqref{eq:epd} из
        вариационного принципа, скобка Ли--Пуансона с тождеством
        Якоби, инварианты и оценки, периодичность, явные решения,
        отсутствие сингулярностей \eqref{eq:noblowup}, оценка МакВини
        с оптимальной константой. Всё это подтверждено численно
        (23/23, \S\ref{sec:numerics}).
  \item \emph{Переносится на полное уравнение}: только каскад
        $2D \to 1D \to 0D$ в той части, где он опирается на
        действительные теоремы (Ладыженская, Коул--Хопф); структурная
        аналогия Казимира верна в $2D$, где энстрофия~--- подлинный
        Казимир \cite{arnoldkhesin}.
  \item \emph{Не переносится}: заключение \eqref{eq:noblowup} о
        трёхмерном уравнении. Причина~--- замыкание петли: в полном
        уравнении $S(t) = \operatorname{sym}(\nabla\bu)$ не экзогенен,
        а определяется тем же полем $\bu$ (через уравнение для
        $\bwn$ и оператор Био--Савара). Модель размыкает эту петлю,
        постулируя $S$ заданным, и потому её теорема~\ref{th:strain}
       ~--- условное утверждение вида «\emph{если} штрейн
        интегрируем, \emph{то} сингулярности нет». Именно
        неэкзогенность $S$ делает проблему тысячелетия открытой
        \cite{fefferman,ckn,bkm}; сведение её к конечной размерности
        само по себе не даёт контроля $\|S(t)\|_{L^\infty}$.
  \item \emph{Технические оговорки}: компактность $\RP^3$
        ограничивает состояние, но не скорости (в модели скорости
        ограничены априорной оценкой, в полном уравнении~--- нет);
        энергетическая монотонность требует \textup{(H2)}
        (замечание~\ref{rem:commutator}); $\bbP(\Om)$ не определена в
        нуле, куда траектория сходится при $\bbL \succ 0$.
\end{enumerate}

Итог: модель $P^3 \to \so$~--- строго верифицируемый конечномерный
каркас геометрии турбулентной завихренности, точный в $2D$-пределе и
условно точный в $3D$; его главный вывод для полного уравнения~---
не «сингулярности нет», а «сингулярность эквивалентна
неинтегрируемости штрейна», что совпадает с направлением классических
критериев \cite{bkm,ckn}.

\needspace{0.25\textheight}
\section{Численная верификация}\label{sec:numerics}

\subsection{Протокол}

Все проверки выполняются интегратором Рунге--Кутты четвёртого порядка
с постоянным шагом $h$ в двойной точности:
\begin{gather*}
  y_{k+1} = y_k + \frac{h}{6}\bigl(k_1 + 2k_2 + 2k_3 + k_4\bigr),\\
  \begin{aligned}
    k_1 &= f(t_k, y_k), &\quad k_2 &= f\bigl(t_k + \tfrac h2,\; y_k + \tfrac h2 k_1\bigr),\\
    k_3 &= f\bigl(t_k + \tfrac h2,\; y_k + \tfrac h2 k_2\bigr), &\quad k_4 &= f(t_k + h,\; y_k + h k_3),
  \end{aligned}
\end{gather*}
порядок сходимости контролируется удвоением шага (Ричардсон):
$e(h)/e(h/2) = 2^{p}$, $p \approx 4$. Исполнение организовано через
суверенный движок \texttt{poler-engine~0.61.0} (режим
\texttt{--exec}): скрипт верификации запускается движком как единая
команда, а скалярные кросс-проверки ниже в подразделе~\ref{sec:calc}
выполнены его встроенным калькулятором (решатель полиномов
Дюрана--Кернера). Среда: numpy~2.1.3, детерминированный датчик
(seed $= 20261002$), полное время прогона $47$~с.

Символьные параметры документа конкретизируются так:
$\bbI = \diag(0{,}90;\,0{,}62;\,0{,}37)$ (общее положение),
$\Om_0 = (2{,}2;\,-1{,}1;\,0{,}7)$; для диссипативных тестов
$\nu = 0{,}10$, $\bbL = \diag(1;\,1{,}7;\,2{,}3)$; осесимметрия
$I_1 = I_2 = 1$, $I_3 = 0{,}5$, $\bbL = \diag(1;1;2)$, $\nu = 0{,}15$;
контрпример замечания~\ref{rem:commutator}~---
$\bbI^{-1} = \diag(10;1;1)$,
$\bbL = \bigl(\begin{psmallmatrix}1&1&0\\1&1{,}5&0\\0&0&1\end{psmallmatrix}\bigr)$;
Бюргерс: $\nu = 0{,}05$ и $0{,}10$, $u_0 = \sin x$,
$\phi_0 = \exp\bigl(-(1-\cos x)/(2\nu)\bigr)$, $N = 2048$ гармоник,
$\phi$ эволюционирует спектрально точно.

\subsection{Результаты}

\begin{longtable}{@{}l >{\raggedright\arraybackslash}p{6.8cm} l l@{}}
\caption{Верификация модели: 23 проверки, все пройдены. Допуски указаны
в описаниях; шаги и горизонты~--- в протоколе и описаниях.}\label{tab:verify}\\
\toprule
ИД & Проверка (допуск) & Измерено & Статус \\
\midrule
\endfirsthead
\caption[]{(продолжение таблицы~\ref{tab:verify})}\\
\toprule
ИД & Проверка (допуск) & Измерено & Статус \\
\midrule
\endhead
\bottomrule
\endfoot
C1 & Дрейф Казимира $C=\|\Om\|^2$, $\nu=0$, $h=10^{-2}$, $T=200$ \ ($\le 2\cdot10^{-8}$, $p\in[3{,}5;5{,}5]$) & $5{,}73\cdot10^{-9}$, $p=4{,}96$ & PASS \\
C2 & Дрейф гамильтониана $\mathcal H$, $\nu=0$ \ ($\le 2\cdot10^{-8}$) & $7{,}50\cdot10^{-9}$ & PASS \\
P1 & Замкнутость орбиты: $d(T_p)/\|\Om_0\|$ \ ($\le 10^{-5}$) & $4{,}98\cdot10^{-15}$, $T_p=3{,}2603$ & PASS \\
D1 & Невязка $\dot C = -2\nu\langle\Om,\bbL\Om\rangle$ \ ($\le 5\cdot10^{-4}$) & $7{,}75\cdot10^{-5}$ & PASS \\
D2 & Монотонность $C(t)$ при $\nu>0$ \ ($\le 0$) & $0{,}00$ & PASS \\
D3 & $C(t) \le C_0e^{-2\nu\lambda_{\min}t}$ \ ($\le 1+10^{-7}$) & $1{,}000000$ & PASS \\
D4 & Монотонность $\mathcal H$ при \textup{(H2)} \ ($\le 0$) & $0{,}00$ & PASS \\
E1a & Контрпример \ref{rem:commutator}: $\lambda_{\min}(\bbL)$ \ ($>0$) & $0{,}2192$ & PASS \\
E1b & $\dot{\mathcal H}(0) = -\nu\langle\Om,\bbI^{-1}\bbL\Om\rangle$ \ ($=+0{,}15$) & $+0{,}1500$ & PASS \\
E1c & Рост $\mathcal H$ при строго убывающем $C$ & $\mathcal H_{\max}/\mathcal H_0 = 1{,}0102$ & PASS \\
B1 & $\sup\|\Om(t)\|/\|\Om_0\|$, $\nu=0$, $T=1000$ \ ($\le 1+10^{-7}$) & $1{,}0000000002$ & PASS \\
B2 & $\|\Om(t)\| \le \|\Om_0\|e^{-\nu\lambda_{\min}t}$, $T=60$ \ ($\le 1+10^{-9}$) & $1{,}0000000000$ & PASS \\
B2b & $\|\Om(T)\|/\|\Om_0\|$, $T=60$ \ ($\le e^{-6}\approx 2{,}48\cdot10^{-3}$) & $2{,}15\cdot10^{-3}$ & PASS \\
A1 & Явное решение \eqref{eq:exact-nu0}, $h=0{,}25$ \ ($p=4\pm0{,}3$) & $5{,}64\cdot10^{-3}$, $p=3{,}97$ & PASS \\
A2 & Явное решение \eqref{eq:exact-nup}, $\nu=0{,}15$ \ ($p=4\pm0{,}3$) & $2{,}40\cdot10^{-18}$, $p=3{,}99$ & PASS \\
M1 & Точный проектор: $\|\bbP^2-\bbP\|$, $\|\mathcal Q(\bbP)-\bbP\|$ \ ($\le 10^{-13}$) & $2{,}7\cdot10^{-16}$ / $8{,}9\cdot10^{-16}$ & PASS \\
M2 & Острота \eqref{eq:q-bound}: $\max\, \delta_Q/(4\delta_P^2)$, $10^4$ возмущённых $\bbP$ \ ($\le 1$) & $0{,}999997$ & PASS \\
H1a & Невязка Бюргерса, $\nu=0{,}1$: абс./отн. \ (отн. $\le 10^{-3}$) & $8{,}92\cdot10^{-6}$ / $6{,}55\cdot10^{-6}$ & PASS \\
H1b & $\max|u_x|$ при $t=0;\,0{,}5;\,1;\,2;\,5$ \ ($<10$; убывание после $t\approx1$) & $1{,}00/1{,}83/4{,}43/7{,}96/2{,}48$ & PASS \\
S1a & Интегрируемый штрейн $s_0e^{-t}$: $\max\|\Om\|/\mathrm{bound}$ \ ($\le 1+10^{-9}$) & $1{,}0000000000$ & PASS \\
S1b & Изотропный $S=s_0\mathrm{Id}$: $\mathrm{ratio}(T)$, рост $2{,}56 \to 139{,}63$ = границе \ ($=1$) & $0{,}99999992$ & PASS \\
L1a & $2D$-сечение: $\max|\Omega_3(t)|$ \ ($\le 10^{-12}$) & $0{,}00$ & PASS \\
L1b & Невязка $\dot C_2 = -2\nu L_{11}(\Omega_1^2+\Omega_2^2)$ \ ($\le 5\cdot10^{-4}$) & $1{,}00\cdot10^{-4}$ & PASS \\
\end{longtable}

\subsection{Скалярные кросс-проверки калькулятором движка}\label{sec:calc}

Два ключевых скалярных факта раздела~\ref{sec:cascade} проверены
встроенным калькулятором \texttt{poler-engine} независимо от
Python-конвейера:
\[
  \texttt{calc\; 3*0.5\^{}2 - 2*0.5\^{}3} \;\longrightarrow\; 0{,}5
  \qquad\text{(неподвижная точка } q(\tfrac12)=\tfrac12\text{),}
\]
\[
  \texttt{calc\; 0.25\^{}2*(3-1)*(2*0.5+1)} \;\longrightarrow\; 0{,}25
  \;=\; 4\cdot(0{,}25)^2
  \qquad\text{(равенство в \eqref{eq:q-bound} при } \lambda=\tfrac12\text{).}
\]
Вторая проверка~--- числовое подтверждение оптимальности константы $4$
(теорема~\ref{th:mcweeny}): левая и правая части оценки совпадают.

\subsection{Комментарии к таблице}

\emph{Инварианты (C1--C2, P1).} Дрейф Казимира и гамильтониана при
$\nu = 0$ имеет порядок $O(h^4)$: измеренный порядок $p = 4{,}96$
подтверждает асимптотику RK4, а замкнутость орбиты
$d(T_p)/\|\Om_0\| = 5\cdot10^{-15}$~--- машинная точность~---
реализует утверждение предложения~\ref{prop:periodic}: численная
траектория возвращается в начальную точку с погрешностью на уровне
двойной точности.

\emph{Диссипация и тонкость (D1--E1).} Тождество
\eqref{eq:casimir} выполняется вдоль траектории с невязкой
$7{,}8\cdot10^{-5}$ (предел конечно-разностной проверки), монотонность
$C$ и $\mathcal H$~--- с нулевым превышением. Контрпример
замечания~\ref{rem:commutator} воспроизведён точно: аналитическое
$\dot{\mathcal H}(0) = +0{,}15$ совпадает с расчётным до $10^{-12}$,
и гамильтониан вырастает на $1{,}02\%$ при строго убывающем Казимире~---
коммутационная гипотеза \textup{(H2)} действительно необходима.

\emph{МакВини и Хопф--Коул (M1--H1).} Точный ранг-1 проектор модели
идемпотентен до $2{,}7\cdot10^{-16}$, а максимальное отношение
$\delta_Q/(4\delta_P^2) = 0{,}999997$ по $10^4$ случайных возмущённых
проекторов (спектр в максимуме $\{0;\,0{,}5017;\,1\}$) подтверждает
и справедливость, и остроту оценки \eqref{eq:q-bound}. Невязка
тождества Коула--Хопфа при $\nu = 0{,}1$ составляет $6{,}6\cdot10^{-6}$
относительно величин слагаемых (усиление округления через $1/\phi$ с
динамическим диапазоном $e^{1/(2\nu)} = e^5$); поведение
$\max|u_x|$ демонстрирует переформирование фронта до $t\approx1$ с
последующим вязким сглаживанием~--- глобальную регулярность
следствия~\ref{cor:burgers-global}.

\emph{Границы метода (S1, B1--B2).} Оценка Грёнуолла
теоремы~\ref{th:strain} выполняется с нулевым нарушением при
интегрируемом штрейне и достигается \emph{точно} при изотропном
$S = s_0\mathrm{Id}$: рост $\|\Om\|$ от $2{,}56$ до $139{,}63$
совпадает с границей во всех знаках~--- константа в \eqref{eq:strain-bound}
неулучшаема, и без условия интегрируемости штрейна экспоненциальный
рост реализуется. Тем самым численно воспроизведена точная граница
метода, установленная в \S\ref{sec:critique}.

Полный код верификации~--- самодостаточный скрипт
\texttt{verify\_\allowbreak so3.py} (Python/numpy, без внешних зависимостей),
поставляемый вместе с документом; трасса прогона~---
\texttt{verification\_\allowbreak output.txt}. Скрипт детерминирован
(фиксированный seed) и воспроизводит таблицу~\ref{tab:verify} одной
командой.


% ============ Литература ============
\begin{thebibliography}{15}
\small
\raggedright

\bibitem{arnold} В.~И.~Арнольд. \emph{Математические методы классической механики}. М.: Наука, 1974.

\bibitem{arnoldkhesin} V.~I.~Arnold, B.~A.~Khesin. \emph{Topological Methods in Hydrodynamics}. Applied Mathematical Sciences 125, Springer, 1998.

\bibitem{marsdenratiu} J.~E.~Marsden, T.~S.~Ratiu. \emph{Introduction to Mechanics and Symmetry}. 2nd ed., Texts in Applied Mathematics 17, Springer, 1999.

\bibitem{hmr} D.~D.~Holm, J.~E.~Marsden, T.~S.~Ratiu. The Euler--Poincar\'e equations and semidirect products. \emph{Advances in Mathematics} 117 (1996), 98--138.

\bibitem{ladyzh} О.~А.~Ладыженская. \emph{Математическая теория вязкой несжимаемой жидкости}. М.: Наука, 1970.

\bibitem{hopf} E.~Hopf. The partial differential equation $u_t+uu_x=\mu u_{xx}$. \emph{Communications on Pure and Applied Mathematics} 3 (1950), 201--230.

\bibitem{cole} J.~D.~Cole. On a quasi-linear parabolic equation occurring in aerodynamics. \emph{Quarterly of Applied Mathematics} 9 (1951), 225--236.

\bibitem{mcweeny} R.~McWeeny. Some recent advances in density matrix theory. \emph{Reviews of Modern Physics} 32 (1960), 335--373.

\bibitem{bkm} J.~T.~Beale, T.~Kato, A.~Majda. Remarks on the breakdown of smooth solutions for the 3-D Euler equations. \emph{Communications in Mathematical Physics} 94 (1984), 61--66.

\bibitem{ckn} L.~Caffarelli, R.~Kohn, L.~Nirenberg. Partial regularity of suitable weak solutions of the Navier--Stokes equations. \emph{Communications on Pure and Applied Mathematics} 35 (1982), 771--831.

\bibitem{fefferman} C.~L.~Fefferman. Existence and smoothness of the Navier--Stokes equation. В кн.: \emph{The Millennium Prize Problems}, AMS, 2006, 57--67.

\bibitem{lasalle} J.~P.~LaSalle. \emph{The Stability of Dynamical Systems}. CBMS-NSF Regional Conference Series 25, SIAM, 1976.

\bibitem{hartman} P.~Hartman. \emph{Ordinary Differential Equations}. 2nd ed., Birkh\"auser, 1982.

\bibitem{euler} L.~Euler. Du mouvement de rotation des corps solides autour d'un axe variable. \emph{M\'emoires de l'Acad\'emie des Sciences de Berlin} 14 (1758), 154--193.

\end{thebibliography}

\end{document}
